\documentclass[10pt,a4paper]{amsart}
\usepackage[margin=19mm]{geometry}
\usepackage{amsmath,amssymb,mathtools}
\usepackage[scr=rsfs,cal=euler]{mathalfa}
\usepackage{xfrac}
\usepackage[unicode,hidelinks,bookmarks=false]{hyperref}
\newcommand{\ie}{i.e.\ }
\newcommand{\cf}{cf.\ }
\newcommand{\resp}{respectively\ }
\newcommand{\Subsection}{\subsection}
\providecommand{\nobreakdash}{\nobreak\hbox{-}}
\setlength{\emergencystretch}{2em}
\multlinegap0pt
\usepackage{amssymb}
\usepackage[shortlabels]{enumitem}
\usepackage{tikz}
\newtheorem{proposition}{Proposition}
\newcommand{\gau}[2]{\left[ \begin{array}{c} #1 \\ #2 \end{array} \right]}
\DeclareMathOperator{\ssum}{sum}
\title{Cyclic Sieving Phenomenon for Independent sets of graphs}
\author{Jacob A White}
\date{Source: arXiv:2605.03083v1 (CC BY 4.0)}
\begin{document}
\maketitle

\noindent\emph{Source and licence.} Cyclic Sieving Phenomenon for Independent sets of graphs. Version arXiv:2605.03083v1 (CC BY 4.0), distributed by arXiv under the Creative Commons Attribution 4.0 International licence, http://creativecommons.org/licenses/by/4.0/, as recorded in the arXiv licence metadata for this exact version and shown on the abstract page https://arxiv.org/abs/2605.03083v1. Full licence notice: \texttt{licenses/arxiv-2605.03083-CC-BY-4.0.txt}.\par\smallskip

\noindent\textit{Editorial scope.} Counted unit: the article's proposition prop:closed (source0.tex lines 419--428). The article's complete original proof of it is delivered with it (source0.tex lines 429--462, one \texttt{proof} environment of 34 lines). No mathematics is written, repaired or re-derived: the statement and the proof are the article's own lines, in the article's own notation, and no retained line is substituted or re-flowed.  No further source-local context is needed: every cross-reference of every retained line resolves inside the two delivered spans.  Nothing was omitted from any retained span, so the package carries no omission record.  No retained line contains a citation, so the wrapper carries no bibliography and no bibliography entry.  No retained line contains a figure environment, an image-inclusion command or drawing code, so no image is delivered and the package declares no asset dependency.  Editorial formatting only, in addition: an AMS \texttt{amsart} preamble that restates the article's own declarations and macros used by the retained lines, namely the environments \texttt{proposition} on separate counters, together with the standard packages those lines need.  The retained environments print in this document as Proposition 1 in order of appearance, whereas the article prints the same items with the numbering of its own full document.  The complete native source of the article as distributed in the arXiv e-print for this exact version is bundled unchanged as \texttt{sources/199769/source0.tex}.
\begin{proposition}
Let $r \leq n$ be positive integers, and let $k$ be a positive integer such that $k \leq \frac{n}{r+1}.$
Then the following identities hold:
\begin{enumerate}
\item $i_k(P^r_n,q) = \gau{n-rk+r}{k}$,
\item $i_k(C_n^r, q) = \frac{[rk]}{[k]}i_{k-1}(P_{n-1-2r}^r,q)+q^{rk}i_k(P^r_{n-r},q)$ 
\item $i_k(C_n^r,q) = \frac{[n]}{[n-rk]}\gau{n-rk}{k}.$
\label{prop:closed}
\end{enumerate}
\end{proposition}

\begin{proof}
 We prove the first result by induction on $r$. 
 When $r=0$, then \[i_k(P_n^0,q) = \sum_{A \subset [0,n]: |A| = k} q^{\ssum_0(A)} = \gau{n}{k}. \]

For $r > 1$, let $A \in \mathcal{I}_k(P^r_n).$ Write $A = \{a_1, \ldots, a_k\}$ where $0 \leq a_1 < a_2 < \cdots < a_k \leq n-1.$ 
Define $A' = \{a_i - i+1: i \in [k] \}$. We observe that $A' \in \mathcal{I}_{k}(P_{n-k+1}^{r-1})$, and moreover $\ssum_r(A) = \ssum_{r-1}(A').$ Hence \[i_k(P_n^r,q) = i_k(P_{n-k+1}^{r-1},q) = \gau{(n-k+1)-(r-1)k+(r-1)}{k} = \gau{n-rk+r}{k}\] 
where the middle equality follows from induction. 

Now we prove our second identity.
Let $r < n$ be positive integers, and let $k$ be a positive integer such that $k(r+1) \leq n$. 
For a given $j \in [0,r-1]$, define
\[\mathcal{A}_j = \{A \in \mathcal{I}_k(C_n^r): j \in A \} \]
Define $\varphi_j: \mathcal{A}_j \to \mathcal{I}_{k-1}(P_{n-2r-1}^r)$ by $\varphi_j(A) = \{a - j-r-1: a \in A \setminus \{j\} \}$. Then $\varphi_j$ is a bijection, and we see that 
\[\ssum_r(A) = j+(j+r+1)(k-1) - (r+1)(k-1) + \ssum_r(\varphi_j(A)) = jk + \ssum_r(\varphi_j(A)).\] 
Thus \[i(\mathcal{A}_j,q) = q^{jk} i_{k-1}(P_{n-2r-1}^r,q) .\]

We let $\mathcal{B} = \{A \in \mathcal{I}_k(C_n^r): A \cap [0,r-1] = \emptyset \}$. We define $\theta:\mathcal{B} \to \mathcal{I}_k(P_{n-r}^r)$ by 
$\theta(B) = \{b - r: b \in B \}$. Then $\theta$ is a bijection, and for $B \in \mathcal{B}$ we have $\ssum_r(B) = rk + \ssum_r(\theta(B))$.
Thus, \[i(\mathcal{B},q) = q^{rk} i_k(P_{n-r}^r,q) .\]

Then we see that 
\begin{align*} i_k(C_n^r,q)  & =  \sum_{j=0}^{r-1} i(\mathcal{A}_j,q)  +  i(\mathcal{B},q) 
 =  \sum_{j=0}^{r-1} q^{jk} i_{k-1}(P_{n-2r-1}^r,q)  +  q^{rk}i_k(P_{n-r}^r,q) \\ 
 & =  i_{k-1}(P_{n-2r-1}^r,q)\sum_{j=0}^{r-1} q^{jk}  +  q^{rk}i_k(P_{n-r}^r,q)  
 =  \frac{[rk]}{[k]} i_{k-1}(P_{n-2r-1}^r,q)  +  q^{rk}i_k(P_{n-r}^r,q). \end{align*}

The third identity follows from the first two, via
\begin{align*}
i_k(C_n^r, q) & =  \frac{[rk]}{[k]} i_{k-1}(P_{n-2r-1}^r,q)  +  q^{rk}i_k(P_{n-r}^r,q) \\
 & =  \frac{[rk]}{[k]}\gau{n-2r-1-r(k-1)+r}{k-1}  +  q^{rk}\gau{n-r-rk+r}{k} \\
& =  \frac{[rk]}{[n-rk]}\gau{n-rk}{k}  +  q^{rk}\frac{[n-rk]}{[n-rk]}\gau{n-rk}{k}  \\
 & =  \frac{[rk]+q^{rk}[n-rk]}{[n-rk]}\gau{n-rk}{k}  =   \frac{[n]}{[n-rk]}\gau{n-rk}{k}.
\end{align*}
\end{proof}

\end{document}
